\section*{Overview}

A modular inverse is the number that undoes multiplication modulo \(m\). It appears everywhere once modular arithmetic
stops being a toy: dividing by factorials, normalizing transitions, solving congruences, and implementing formulas that
look rational but are really modular.

\section*{When to Use It}

Use a modular inverse when:

\begin{itemize}
  \item a formula under modulo arithmetic needs division,
  \item you need to solve \(ax \equiv b \pmod m\),
  \item factorial inverses or combinatorics tables are involved.
\end{itemize}

\section*{Core Idea}

The inverse of \(a\) modulo \(m\) is a value \(a^{-1}\) such that
\[
a \cdot a^{-1} \equiv 1 \pmod m.
\]
It exists iff \(\gcd(a, m) = 1\).

There are two common ways to compute it:
\begin{itemize}
  \item extended Euclid for a general modulus,
  \item fast exponentiation \(a^{m-2}\) when \(m\) is prime.
\end{itemize}

\section*{Key Insight}

``Division modulo \(m\)'' is not automatic. It is multiplication by an inverse, and that only makes sense when the
number is coprime with the modulus.

\section*{Operations / Main Technique}

\begin{itemize}
  \item inverse by extended Euclid for arbitrary coprime \((a, m)\),
  \item inverse by Fermat's little theorem for prime \(m\),
  \item batch inverse tables when many inverses under one prime modulus are needed.
\end{itemize}

\paragraph{Worked example.}
Modulo \(11\), the inverse of \(3\) is \(4\) because \(3 \cdot 4 = 12 \equiv 1 \pmod{11}\).

\section*{Correctness Intuition}

Extended Euclid returns coefficients \(x, y\) such that
\[
ax + my = \gcd(a, m).
\]
When the gcd is \(1\), reducing modulo \(m\) gives \(ax \equiv 1 \pmod m\), so \(x\) is the inverse.

For prime moduli, Fermat gives \(a^{m-1} \equiv 1\), so \(a^{m-2}\) is the inverse.

\section*{Complexity Analysis}

\begin{itemize}
  \item inverse by extended Euclid: \(O(\log m)\),
  \item inverse by fast power: \(O(\log m)\),
  \item build inverse table \(1 \ldots n\) under prime modulus: \(O(n)\).
\end{itemize}

\section*{Implementation}

The code includes both the general and prime-modulus versions. I prefer keeping both because contest moduli vary, and
assuming primality when it is not given is a common bug.

\section*{Common Pitfalls}

\begin{itemize}
  \item Trying to invert a value that is not coprime to the modulus.
  \item Using \(a^{m-2}\) when the modulus is not prime.
  \item Forgetting to normalize the extended-Euclid result into \([0, m-1]\).
  \item Dividing under modulo arithmetic before taking the modulo safely.
\end{itemize}

\section*{Variants / Extensions}

\begin{itemize}
  \item Precompute factorials and inverse factorials.
  \item Solve linear congruences and CRT systems.
  \item Inverse tables for repeated combinatorics queries.
\end{itemize}

\section*{Practice Problems}

\begin{itemize}
  \item Compute binomial coefficients modulo a prime.
  \item Solve \(ax \equiv b \pmod m\).
  \item Evaluate rational-looking formulas modulo a large prime.
\end{itemize}

\section*{References}

\begin{itemize}
  \item \href{https://cp-algorithms.com/algebra/module-inverse.html}{cp-algorithms: Modular Multiplicative Inverse}.
  \item \href{https://usaco.guide/gold/modular?lang=cpp}{USACO Guide: Modular Arithmetic}.
  \item \href{https://codeforces.com/catalog}{Codeforces Catalog}.
\end{itemize}
